\chapter[A recapitalization and its owners]{A recapitalization: cash raised, rights transferred, and incentives changed}
\label{sim:ownership-case}

A financing plan can save a company while reducing the value held by its existing owners. It can also dilute their percentage without reducing their value. Distinguishing these outcomes requires a security-level calculation. This chapter follows a synthetic recapitalization with ordinary shares, a preference offer and a convertible claim. The terms are stipulated so the arithmetic can be reconstructed. They illustrate mechanisms rather than legal or valuation conclusions for a real issuer.

\section{The company before financing}
There are 100 ordinary shares. Existing operating assets have a stipulated economic value of 600, cash is 20 and a senior debt payment of 80 is immediately due. For the first comparison, assume asset value is unchanged by financing and there are no other claims. Before paying debt, residual economic equity value is $600+20-80=540$, or 5.4 per ordinary share.

The company nevertheless cannot pay 80 from cash 20. It needs 60 of additional cash. Positive economic equity value does not eliminate this liquidity problem. Selling assets, borrowing, issuing equity or negotiating deferral are different possible responses. We first study an ordinary-share offer.

Book assets may have a different carrying amount. Nothing in the calculation requires book equity to equal 540. If the model uses book equity to set a subscription price of 5.4, it would be making a separate valuation assumption. Here 540 is a supplied economic valuation used only for the illustrative allocation.

\section{A fairly priced ordinary issue}
Suppose investors subscribe at 5.4 with no fees, and the company raises exactly 60. New shares are $60/5.4=100/9$, about 11.1111. Cash becomes 80 and is paid to the senior creditor. Remaining assets are 600, there is no debt, and shares total $1000/9$.

Post-transaction value per share is
\begin{equation}
 \frac{600}{1000/9}=5.4.
\end{equation}
Existing holders retain value 540, and new holders own value 60. Their percentages change, but no value is transferred under these assumptions. The creditor receives its contractual 80. The cash injection solves liquidity without creating extra value by itself.

This result relies on the supplied fair price, homogeneous rights, no costs and unchanged operating value. It is a useful benchmark because deviations can be attributed to specific mechanisms. It is not a prediction that a distressed company can always issue at fair value.

\section{A discounted issue}
Now suppose the offer price is 3, with the same 60 cash target and no fees. Twenty shares are issued. After paying debt, value remains 600 but shares total 120. Each share is worth 5 under the stipulated valuation. Existing holders own 500, a loss of 40 relative to the benchmark; new investors paid 60 for claims worth 100.

The transfer is not an additional operating loss. It is an allocation effect of selling claims below the assumed post-transaction value. A journal recording cash and equity does not reveal it. The ownership register and valuation assumptions are needed.

If investors' required price of 3 reflects information that operating assets are worth less than 600, the benchmark valuation is wrong. The calculation should not accuse the offer of unfairness without validating the valuation premise. Scenario analysis can show transfers conditional on alternative asset values and prices.

\section{Fees change both quantity and value}
Add a fixed cash issuance fee of 6. To net 60, the company must raise 66. At price 3, it issues 22 shares. After fees and debt payment, operating assets remain 600, cash is zero and shares total 122. Existing owners' value is $600(100/122)$, about 491.8033.

The fee destroys 6 of total value relative to an otherwise identical financing if it buys no additional benefit in this model. The discounted price transfers value between old and new owners. These effects should be separated. Calling both dilution obscures which mechanism could be changed by renegotiation or a different financing channel.

If the fee is paid later, current funding need differs and a payable remains. If it is financed by issuing additional shares to an intermediary, share count changes without the same cash payment. The offer must state the form and timing of compensation.

\section{An authorization constraint}
Suppose only fifteen new ordinary shares may be issued under the admitted authorization. At price 3 and fee 6, maximum net cash is 39. Existing cash 20 then supplies only 59 toward debt 80, leaving a shortfall of 21. The ordinary issue cannot solve the maturity under these terms.

At price $p$, feasibility requires $15p-6\geq60$, so $p\geq4.4$. This is an exact offer threshold. It does not establish that investors will pay 4.4. A market model or a new offer is needed to answer that question.

An optimizer that increases shares beyond fifteen has changed the feasible set. An accounting routine that credits capital for the missing 21 without issuing a security has invented financing. A correct result records the shortfall and considers other admitted actions.

\section{A rights offer}
Suppose existing holders have a right to subscribe proportionally at the discounted price. If every holder subscribes in proportion, their relative ownership remains unchanged, but they contribute new cash. The issue can still have fees and execution risk.

A holder who does not subscribe may sell the right if trading is permitted and a buyer exists. The value of the right depends on issue terms and market value. A model that treats all existing holders as equally able and willing to contribute assumes away liquidity and participation differences.

For AccDSL, a rights offer needs entitlement, subscription, transfer and expiry events. Unexercised rights may lapse or be allocated under a backstop arrangement. The company receives cash only from settled subscriptions or the backstop, not from the mere existence of rights. Historical statements may not reveal the eventual participation pattern.

\section{A preference alternative}
Consider instead an investor providing 60 for a nonparticipating preference claim. At a future exit, the investor may take 90 before ordinary shareholders or convert into 20 per cent of ordinary equity. Assume no senior debt remains, no interim dividends, no other rights and distributable value $V$. The investor receives
\begin{equation}
 P(V)=\max\{\min(90,V), 0.2V\}.
 \label{sim:preference-payoff}
\end{equation}
Ordinary holders receive $V-P(V)$. The first term respects limited available value: a preference of 90 cannot pay 90 if only 50 exists. Conversion dominates the full preference when $0.2V>90$, or $V>450$.

For $V=50$, the investor takes 50 and ordinary holders receive zero. For $V=300$, the investor takes 90 and ordinary holders receive 210. For $V=600$, conversion pays 120 and ordinary holders receive 480. The same 60 of initial funding creates very different allocation across outcomes from a simple ordinary issue.

\section{Participation changes the waterfall}
A participating preference can take a preference amount and then share in remaining value. Suppose the investor receives up to 90 first and then 20 per cent of the residual, with no cap. Its payoff is
\begin{equation}
 P_{\rm part}(V)=\min(90,V)+0.2\max(V-90,0).
\end{equation}
At $V=300$, this gives $90+42=132$, compared with 90 for the nonparticipating offer. At $V=600$, it gives 192, compared with 120 under conversion in the earlier contract.

These are different securities, not two calculation methods for one security. A parser that extracts only the word preference and the amount 90 will miss the distinction. The contract admission model must preserve participation, conversion and caps where they affect the query.

If a participation cap applies, another region boundary appears. The engine can represent the payoff as a piecewise function and split only where the active right changes. A sample of a few exit values is useful for illustration but not a complete characterization of the payoff.

\section{Value conservation in each state}
For every exit value, distributions to all claimants must not exceed distributable assets under the assumed waterfall. In the nonparticipating example, $P(V)+[V-P(V)]=V$ exactly. This is a conservation check independent of whether the offer is fairly priced.

A common implementation error pays the preference and also allocates 20 per cent of total value as though conversion were simultaneous. That double counts mutually exclusive rights. Another pays 90 when $V<90$, creating cash. Both can be detected with low-value boundary cases.

Fair pricing requires discounting and a distribution or valuation operator over future payoffs. The waterfall itself supplies state-contingent claims, not their current price. An investor's willingness to pay 60 can be taken as a declared offer without claiming that a model has proved its fairness.

\section{A convertible debt alternative}
Suppose an existing note has face 80 due at exit and can convert into sixteen ordinary shares, with no other debt. If conversion is optional for the holder and the company has 100 ordinary shares before conversion, conversion gives a fraction $16/116$ of total equity after conversion. The holder compares that claim with the debt payoff under the specified default and priority rules.

If distributable asset value is high enough that debt is paid in full, conversion value is $16V/116$. It exceeds 80 when $V>580$. Below that threshold, the holder may prefer debt. At low asset values, debt recovery is limited by available assets and seniority. The conversion threshold can change if other securities exist.

A mistaken calculation using $16V/100$ ignores the new shares in the denominator and gives a threshold of 500. That is a dilution error. A second mistake pays debt 80 and also issues sixteen shares. The conversion event must extinguish the specified debt claim if that is what the contract says.

\section{Cash settlement and issuer options}
Some instruments permit cash settlement of conversion value or allow the issuer to choose a settlement form. That changes liquidity and ownership. A model that always assumes physical shares may understate cash need; one that always assumes cash may overstate dilution avoidance.

The choice belongs to the contract and the authorized actor. If the issuer can choose, its policy must consider cash capacity and economic cost. If the holder chooses, management cannot select the favorable form merely to balance the forecast.

An unknown settlement clause should produce conditional outcomes. The system can show a physical-settlement branch and a cash-settlement branch with their premises. It should seek the clause when the distinction changes a liquidity conclusion. This is a high-value source-research target.

\section{Control and approvals}
A recapitalization can shift voting control even when economic ownership changes modestly. A preference investor may receive veto rights or board appointment rights. Those rights can affect future financing, dividends and asset sales. Treating them as irrelevant because they have no immediate journal amount can distort later policies.

The model need not simulate every governance detail for every query. It should retain material rights and declare when behavior assumes their exercise or nonexercise. If a future equity issue requires investor consent, the feasible action set must include that condition.

A public-data analyst may not know whether consent will be given. This is future behavioral uncertainty, distinct from uncertainty about whether the consent right exists. Combining both into one unspecified branch makes it harder to identify useful evidence or policy assumptions.

\section{Incentives after financing}
Suppose ordinary holders choose between a safe project with exit value 300 and a risky project yielding 100 or 600 with equal probability. Under the nonparticipating preference, ordinary payoff from the safe project is 210. Under the risky project, ordinary payoffs are 10 and 480, with mean 245.

Total expected exit value of the risky project is 350, above 300, so both total value and ordinary value favor it in this particular example. Change the low outcome to zero and high outcome to 500. Total expected value is 250, below 300, while ordinary expected payoff is $[0+400]/2=200$, also below 210. The rankings still agree. These examples should not be used to claim a general conflict.

To exhibit a conflict, consider a senior fixed claim of 250 and ordinary residual value. A safe project with value 300 gives ordinary holders 50. A risky project with values zero and 500 equally likely has total expected value 250 but ordinary expected payoff 125. Ordinary holders prefer the risky project while total expected value is lower. The contract allocation creates the divergence under the stated objective and probabilities.

\section{Why the actor objective matters}
A policy maximizing total enterprise value selects the safe project in the final example. A policy maximizing ordinary equity selects the risky one. A lender pricing debt must anticipate the policy actually modeled. Pricing under the safe policy and then letting shareholders choose the risky one is inconsistent.

Covenants or control rights may constrain the choice. Their presence changes the feasible set and possibly the price offered initially. This joins the security register to behavioral modeling. The register is not merely a reporting table; it determines who receives gains and losses and who can authorize actions.

The example is a local derivation of incentive conflict, not a calibrated prediction. It gives a discriminating fixture for a structural solver: the actor objective, priority and feasible actions should reproduce the stated ranking. A solver returning the same policy for both objectives without a binding constraint is suspect.

\section{Book entries and economic allocation}
The cash issue creates cash and accounting equity under the applicable treatment. The preference or convertible may require classification and component allocation that cannot be inferred from the simplified payoff alone. Those accounting decisions need the standard, jurisdiction and complete terms.

Economic valuation can proceed conditionally on a specified payoff while accounting classification remains unresolved. Conversely, a classified instrument can still have uncertain market price. The system should allow these statuses to coexist. Forcing both layers into one known or unknown flag loses useful work.

A book-equity bridge should reconcile issue proceeds, fees under the chosen treatment, retained earnings, distributions and other movements. A separate ownership bridge reconciles units and rights. A valuation bridge explains changes in economic claims. The three bridges communicate but should not be substituted for one another.

\section{Historical ambiguity in a capital movement}
Suppose a filing reports net cash from equity transactions of 60 and an increase of twenty ordinary shares. This could be a simple issue at 3 with no costs. It could also be an issue of thirty shares at 4 and repurchase of ten at 6, again netting to 60 and twenty shares. The closing share count and net cash agree.

The histories differ in gross funding, repurchase behavior and prices. If future policy extrapolates issuance discounts or repurchase commitments, the distinction matters. A note disclosing gross proceeds and repurchases can narrow the family. Without it, a single implied price of 3 is not established as an actual transaction price.

This example parallels the receivables and debt bridges. Net movements constrain sums but do not identify gross events. The security register should therefore permit unresolved transaction families while preserving observed opening and closing counts.

\section{An end-to-end fixture}
A useful fixture starts with the 100-share company, the debt maturity and the cash need. It evaluates ordinary issuance at fair and discounted prices, fixed fees, authorization limits, preference rights and conversion. Each transaction updates cash, liabilities, shares and rights before producing statements.

Expected failures include using book value as an unstated subscription price, ignoring the authorization cap, counting conversion shares against the old denominator, paying mutually exclusive preference and conversion rights together, and treating an unsigned offer as settled cash. Each failure has a specific input that exposes it.

The fixture should also retain unresolved accounting classification as a legitimate result where full terms are absent. A financing calculation can be mechanically complete while a statement presentation remains conditional. That distinction is preferable to inventing a universal accounting rule to make the example look finished.

\section{What the case adds to simulation design}
Equity funding is a transaction among actors with claims, rights and constraints. Its effect cannot be reduced to an external-equity cost scalar when ownership and future decisions matter. Nor can share count alone establish value transfer. Prices, fees, rights and valuation premises jointly determine the result.

The proposed AccDSL design therefore keeps a dated security register, offer and settlement states, class-specific payoff rules and actor objectives. It can use a simple homogeneous-share profile when appropriate, but retains the boundary that tells a user when richer terms are required. This makes dilution and financing consequences explicit without pretending that every security must always be simulated at maximum detail.
